\documentclass[11pt]{article}

\usepackage[margin=1in]{geometry}
\usepackage{amsmath,amssymb,amsthm}
\usepackage{booktabs}
\usepackage{longtable}
\usepackage{array}
\usepackage{enumitem}
\usepackage{xcolor}
\usepackage{hyperref}
\usepackage{titlesec}

\hypersetup{
    colorlinks=true,
    linkcolor=blue,
    urlcolor=blue,
    citecolor=blue
}

\setlist[itemize]{leftmargin=1.5em}
\setlist[enumerate]{leftmargin=1.8em}

\title{\textbf{AS.110.106 Calculus I (Biology and Social Sciences)}\\
\large Full Course Breakdown and Learning Framework}
\author{Department of Mathematics --- Johns Hopkins University}
\date{}

\begin{document}

\maketitle

\section{Course Overview}

\subsection{Course Identity}

\begin{itemize}
    \item \textbf{Course:} AS.110.106 Calculus I (Biology and Social Sciences)
    \item \textbf{Discipline:} Mathematics / Single-Variable Calculus
    \item \textbf{Primary audience:} Students in biology, medicine-related fields, social sciences, and other quantitative disciplines.
    \item \textbf{Course emphasis:} Functions, limits, derivatives, applications of derivatives, integration, and techniques/applications of integration.
    \item \textbf{Estimated core duration:} Approximately 12--13+ weeks.
\end{itemize}

\subsection{Primary Text}

\begin{quote}
C. Neuhauser and M. Roper, \textit{Calculus for Biology and Medicine}, 4th Edition, Pearson, 2018.\\
ISBN-10: 0134070046 \qquad ISBN-13: 978-0134070049
\end{quote}

\subsection{Course Purpose}

This course develops the central ideas and computational tools of single-variable calculus with emphasis on interpretation and applications in biology, medicine, population dynamics, and social sciences. Students learn to model changing quantities, determine instantaneous and average rates of change, optimize outcomes under mathematical constraints, and calculate accumulated change using definite integrals.

\section{Prerequisite Knowledge}

Students should enter the course with fluency in algebra, functions, graphs, and basic trigonometry. The course assumes these skills and uses them continuously.

\subsection{Essential Prerequisites}

\begin{longtable}{p{0.26\textwidth} p{0.67\textwidth}}
\toprule
\textbf{Prerequisite Area} & \textbf{Expected Knowledge and Skills} \\
\midrule
\endhead

Algebra &
Manipulating algebraic expressions; factoring; simplifying rational expressions; solving linear, quadratic, exponential, and logarithmic equations; working with inequalities. \\

Functions &
Function notation; domain and range; composition of functions; inverse functions; transformations; interpreting functions from equations, graphs, tables, and word problems. \\

Graphing and Coordinate Geometry &
Graphing linear, polynomial, rational, exponential, logarithmic, and trigonometric functions; intercepts; asymptotes; slopes; reading graphical trends. \\

Exponents and Logarithms &
Laws of exponents; laws of logarithms; natural logarithm $\ln x$; exponential functions $e^x$; solving exponential and logarithmic equations. \\

Trigonometry &
Radians; unit circle values; sine, cosine, tangent; basic identities; inverse trigonometric functions; graph behavior of trigonometric functions. \\

Quantitative Modeling &
Translating a verbal problem into variables, equations, functions, constraints, and appropriate units. \\

Arithmetic and Numerical Reasoning &
Working accurately with fractions, decimals, percentages, scientific notation, estimation, and calculator-based computation. \\
\bottomrule
\end{longtable}

\subsection{Recommended Readiness Topics}

Students who need additional preparation should review:

\begin{itemize}
    \item Factoring polynomials and solving quadratic equations.
    \item Rational-expression simplification and common denominators.
    \item Function composition, inverse functions, and domain restrictions.
    \item Exponential growth and decay models.
    \item Natural logarithms and the exponential function $e^x$.
    \item Radian measure and basic trigonometric identities.
    \item Interpretation of slope, average rate of change, and graphs.
\end{itemize}

\section{Course-Level Learning Objectives}

By the end of the course, students should be able to:

\begin{enumerate}
    \item Analyze functions using symbolic, graphical, numerical, and verbal representations.
    \item Model exponential growth and decay in biological, medical, and social-science settings.
    \item Evaluate and interpret limits, including limits at infinity and trigonometric limits.
    \item Determine whether a function is continuous and explain the mathematical and applied significance of continuity.
    \item Define the derivative as an instantaneous rate of change and as the slope of a tangent line.
    \item Compute derivatives using the power, product, quotient, chain, implicit differentiation, and inverse-function rules.
    \item Apply derivatives to analyze rates, monotonicity, concavity, extrema, inflection points, and optimization problems.
    \item Construct and interpret linear approximations of nonlinear functions.
    \item Compute indefinite and definite integrals.
    \item Apply the Fundamental Theorem of Calculus to connect differentiation, accumulation, and area.
    \item Use integration to solve accumulation and applied-rate problems.
    \item Select and execute appropriate integration techniques, including substitution, integration by parts, and partial fractions.
    \item Communicate mathematical reasoning clearly, with correct notation, units, assumptions, and interpretation.
\end{enumerate}

\section{Course Structure}

\subsection{Module Sequence}

\begin{longtable}{p{0.08\textwidth} p{0.27\textwidth} p{0.13\textwidth} p{0.42\textwidth}}
\toprule
\textbf{Module} & \textbf{Major Topic} & \textbf{Estimated Time} & \textbf{Core Focus} \\
\midrule
\endhead

1 & Functions, Exponential Models, and Sequences & 2 weeks &
Review mathematical functions; develop exponential growth and decay models; introduce sequences as discrete patterns and precursors to limits. \\

2 & Limits and Continuity & 2 weeks &
Develop the conceptual foundation of calculus through limiting processes, continuity, limits at infinity, and key theorems. \\

3 & Derivatives and Differential Rules & 3 weeks &
Define derivatives formally and develop techniques for differentiating algebraic, trigonometric, exponential, inverse, and implicit functions. \\

4 & Applications of Differentiation & 2+ weeks &
Use derivatives to study shape, behavior, optimization, approximation, and limiting forms. \\

5 & Integration and the Fundamental Theorem & 2 weeks &
Develop antiderivatives, definite integrals, accumulation, and the relationship between differentiation and integration. \\

6 & Techniques and Applications of Integration & 1+ week &
Apply substitution, integration by parts, and partial fractions to evaluate more complex integrals. \\
\bottomrule
\end{longtable}

\section{Detailed Module Breakdown}

\subsection{Module 1: Functions, Exponential Growth, Decay, and Sequences}

\textbf{Estimated duration:} 2 weeks

\subsubsection{Syllabus Chapters and Sections}

\begin{itemize}
    \item Chapter 1: Review of Basic Properties of Functions
    \item Section 2.1: Exponential Growth and Decay
    \item Section 2.2: Sequences
\end{itemize}

\subsubsection{Major Concepts}

\begin{itemize}
    \item Function notation, domain, range, and graphical interpretation.
    \item Function transformations and composition.
    \item Inverse functions and one-to-one behavior.
    \item Average rate of change.
    \item Exponential functions of the form
    \[
        P(t) = P_0 e^{kt}.
    \]
    \item Growth versus decay according to the sign of $k$.
    \item Doubling time and half-life.
    \item Discrete sequences and recursive descriptions.
    \item Long-term behavior of sequences.
\end{itemize}

\subsubsection{Learning Objectives}

Students will be able to:

\begin{enumerate}
    \item Identify domain, range, intercepts, asymptotes, and transformations of common functions.
    \item Evaluate and compose functions using formulas, tables, and graphs.
    \item Interpret average rate of change in context.
    \item Build and solve exponential growth and decay models.
    \item Determine the meaning of parameters such as $P_0$ and $k$ in an exponential model.
    \item Calculate doubling time or half-life using logarithms.
    \item Represent discrete processes using explicit and recursive sequences.
    \item Describe whether a sequence converges, diverges, increases, decreases, or oscillates.
\end{enumerate}

\subsubsection{Applied Contexts}

\begin{itemize}
    \item Population growth.
    \item Spread of organisms or cells under simplified assumptions.
    \item Radioactive decay and medical tracer decay.
    \item Drug elimination models.
    \item Compound interest and economic growth.
    \item Discrete-time population or resource models.
\end{itemize}

\subsubsection{Difficult Topics}

\begin{itemize}
    \item Distinguishing exponential from linear growth using data or graphs.
    \item Solving for time in equations involving exponentials.
    \item Correctly interpreting a negative exponential rate constant.
    \item Determining realistic domains in applied models.
    \item Connecting discrete sequences to continuous functions and limits.
\end{itemize}

\subsection{Module 2: Limits and Continuity}

\textbf{Estimated duration:} 2 weeks

\subsubsection{Syllabus Chapters and Sections}

\begin{itemize}
    \item Section 3.1: Limits
    \item Section 3.2: Continuity
    \item Section 3.3: Limits at Infinity
    \item Section 3.4: Trigonometric Limits and the Sandwich Theorem
    \item Section 3.5: Properties of Continuous Functions
    \item Section 3.6: Formal Definitions of Limits and Limits at Infinity
\end{itemize}

\subsubsection{Major Concepts}

\begin{itemize}
    \item Informal and formal definitions of limits.
    \item One-sided limits.
    \item Infinite limits and vertical asymptotes.
    \item Limits at infinity and horizontal asymptotes.
    \item Continuity at a point and on an interval.
    \item Removable, jump, and infinite discontinuities.
    \item The Sandwich Theorem.
    \item The fundamental trigonometric limit
    \[
        \lim_{x \to 0}\frac{\sin x}{x} = 1.
    \]
    \item Intermediate Value Theorem and other properties of continuous functions.
\end{itemize}

\subsubsection{Learning Objectives}

Students will be able to:

\begin{enumerate}
    \item Estimate limits from graphs, tables, and numerical evidence.
    \item Evaluate limits algebraically using direct substitution, factoring, rationalization, and limit laws.
    \item Determine and interpret one-sided limits.
    \item Analyze infinite limits and limits at infinity.
    \item Identify horizontal and vertical asymptotic behavior.
    \item Determine whether a function is continuous at a point or over an interval.
    \item Classify discontinuities.
    \item Use the Sandwich Theorem to establish selected limits.
    \item Explain the formal meaning of a limit using $\varepsilon$--$\delta$ language.
    \item Apply properties of continuous functions to justify the existence of solutions.
\end{enumerate}

\subsubsection{Difficult Topics}

\begin{itemize}
    \item Understanding that a limit concerns nearby behavior rather than necessarily the function value at a point.
    \item Evaluating indeterminate forms such as $\frac{0}{0}$ without assuming the limit is zero.
    \item Distinguishing divergence from an infinite limit.
    \item Interpreting limits at infinity versus limits as $x$ approaches a finite number.
    \item Applying the formal $\varepsilon$--$\delta$ definition.
    \item Choosing an appropriate comparison for the Sandwich Theorem.
\end{itemize}

\subsubsection{Conceptual Dependency}

\[
\text{Functions} \longrightarrow \text{Limits} \longrightarrow \text{Continuity}
\longrightarrow \text{Derivative Definition}
\]

\subsection{Module 3: Derivatives and Rules of Differentiation}

\textbf{Estimated duration:} 3 weeks

\subsubsection{Syllabus Chapters and Sections}

\begin{itemize}
    \item Section 4.1: Formal Definition of the Derivative
    \item Section 4.2: Basic Rules of Derivatives
    \item Section 4.3: Power Rule and Rules of Differentiation
    \item Section 4.4: Product and Quotient Rules
    \item Section 4.5: Chain Rule
    \item Section 4.6: Implicit Functions and Implicit Differentiation
    \item Section 4.7: Higher Derivatives
    \item Section 4.8: Derivatives of Trigonometric Functions
    \item Section 4.9: Derivatives of Exponential Functions
    \item Section 4.10: Derivatives of Inverse Functions
    \item Section 4.11: Linear Approximation
\end{itemize}

\subsubsection{Major Concepts}

\begin{itemize}
    \item Derivative as a limit of average rates of change:
    \[
        f'(x) = \lim_{h \to 0}\frac{f(x+h)-f(x)}{h}.
    \]
    \item Derivative as slope of a tangent line.
    \item Derivative as instantaneous rate of change.
    \item Power, constant, sum, and difference rules.
    \item Product rule:
    \[
        \frac{d}{dx}[f(x)g(x)] = f'(x)g(x) + f(x)g'(x).
    \]
    \item Quotient rule:
    \[
        \frac{d}{dx}\left[\frac{f(x)}{g(x)}\right]
        =
        \frac{g(x)f'(x)-f(x)g'(x)}{[g(x)]^2}.
    \]
    \item Chain rule:
    \[
        \frac{d}{dx}[f(g(x))] = f'(g(x))g'(x).
    \]
    \item Implicit differentiation.
    \item Higher-order derivatives.
    \item Derivatives of trigonometric, exponential, logarithmic, and inverse functions.
    \item Linearization and differential approximation:
    \[
        L(x) = f(a) + f'(a)(x-a).
    \]
\end{itemize}

\subsubsection{Learning Objectives}

Students will be able to:

\begin{enumerate}
    \item Use the formal limit definition to compute derivatives of elementary functions.
    \item Interpret derivatives graphically, numerically, and in applied settings.
    \item Compute derivatives using basic differentiation rules.
    \item Apply product, quotient, and chain rules accurately.
    \item Differentiate composite functions with multiple nested layers.
    \item Differentiate implicitly defined relationships.
    \item Compute and interpret second and higher derivatives.
    \item Differentiate trigonometric, exponential, logarithmic, and inverse functions.
    \item Find equations of tangent and normal lines.
    \item Use linear approximation to estimate function values near a known point.
    \item Interpret derivatives with appropriate units.
\end{enumerate}

\subsubsection{Applied Contexts}

\begin{itemize}
    \item Instantaneous population growth rates.
    \item Velocity and acceleration.
    \item Marginal cost, marginal revenue, and marginal benefit.
    \item Rates of chemical reaction.
    \item Changes in concentration over time.
    \item Sensitivity of a modeled biological response.
    \item Approximation of nonlinear quantities near baseline conditions.
\end{itemize}

\subsubsection{Difficult Topics}

\begin{itemize}
    \item Computing a derivative directly from the difference quotient.
    \item Recognizing when the chain rule is required.
    \item Correctly applying product, quotient, and chain rules in one problem.
    \item Implicit differentiation, especially when both variables depend on one another.
    \item Differentiating inverse functions and inverse trigonometric functions.
    \item Distinguishing between a function, its derivative, and its second derivative.
    \item Building and interpreting linear approximations with units and error awareness.
\end{itemize}

\subsection{Module 4: Applications of Differentiation}

\textbf{Estimated duration:} 2+ weeks

\subsubsection{Syllabus Chapters and Sections}

\begin{itemize}
    \item Section 5.1: Extreme and Mean Value Theorems
    \item Section 5.2: Monotonicity and Concavity
    \item Section 5.3: Extrema and Inflection Points
    \item Section 5.4: Optimization
    \item Section 5.5: L'Hopital's Rule
\end{itemize}

\subsubsection{Major Concepts}

\begin{itemize}
    \item Extreme Value Theorem.
    \item Mean Value Theorem.
    \item Critical points.
    \item Increasing and decreasing behavior.
    \item Concavity and second derivatives.
    \item Local and absolute extrema.
    \item Inflection points.
    \item First derivative and second derivative tests.
    \item Optimization under constraints.
    \item Indeterminate forms and L'Hopital's Rule.
\end{itemize}

\subsubsection{Learning Objectives}

Students will be able to:

\begin{enumerate}
    \item State and apply the Extreme Value Theorem and Mean Value Theorem under appropriate conditions.
    \item Find critical points of a function.
    \item Determine intervals on which a function is increasing or decreasing.
    \item Determine concavity and locate possible inflection points.
    \item Identify local and absolute maxima and minima.
    \item Use first and second derivative tests to classify critical points.
    \item Sketch or analyze a function using derivative information.
    \item Formulate and solve applied optimization problems.
    \item Evaluate appropriate indeterminate forms using L'Hopital's Rule.
    \item Interpret optimization results in terms of units, feasibility, and the original applied setting.
\end{enumerate}

\subsubsection{Optimization Problem Framework}

A standard optimization solution should include:

\begin{enumerate}
    \item Define variables with units.
    \item Identify the quantity to maximize or minimize.
    \item Write the objective function.
    \item Use the constraint to express the objective in one variable.
    \item State the feasible domain.
    \item Differentiate the objective function.
    \item Find and test critical points and relevant endpoints.
    \item Report the result using a complete contextual sentence.
\end{enumerate}

\subsubsection{Difficult Topics}

\begin{itemize}
    \item Distinguishing critical points from actual extrema.
    \item Performing a full first-derivative sign analysis.
    \item Identifying valid inflection points rather than merely solving $f''(x)=0$.
    \item Constructing an objective function and constraint from words.
    \item Checking domain restrictions and endpoints in optimization.
    \item Recognizing the permitted indeterminate forms for L'Hopital's Rule.
    \item Avoiding misuse of L'Hopital's Rule when algebraic simplification is preferable.
\end{itemize}

\subsubsection{Conceptual Dependency}

\[
\text{Derivative Rules}
\longrightarrow
\text{First and Second Derivatives}
\longrightarrow
\text{Function Analysis}
\longrightarrow
\text{Optimization and Modeling}
\]

\subsection{Module 5: Integration and the Fundamental Theorem of Calculus}

\textbf{Estimated duration:} 2 weeks

\subsubsection{Syllabus Chapters and Sections}

\begin{itemize}
    \item Section 5.10: Antiderivatives
    \item Section 6.1: The Definite Integral
    \item Section 6.2: The Fundamental Theorem of Calculus
    \item Section 6.3: Applications of Integration
\end{itemize}

\subsubsection{Major Concepts}

\begin{itemize}
    \item Antiderivatives and indefinite integrals.
    \item Constants of integration.
    \item Definite integrals as net accumulation.
    \item Area under and between curves.
    \item Riemann sums and partitioning.
    \item Fundamental Theorem of Calculus, Part I:
    \[
        \frac{d}{dx}\int_a^x f(t)\,dt = f(x).
    \]
    \item Fundamental Theorem of Calculus, Part II:
    \[
        \int_a^b f(x)\,dx = F(b)-F(a),
    \]
    where $F'(x)=f(x)$.
    \item Accumulation functions.
    \item Applications involving total change, area, and accumulated quantities.
\end{itemize}

\subsubsection{Learning Objectives}

Students will be able to:

\begin{enumerate}
    \item Find antiderivatives of elementary functions.
    \item Use initial conditions to determine constants of integration.
    \item Interpret a definite integral as a net accumulated change.
    \item Estimate definite integrals from graphs, tables, and Riemann sums.
    \item Evaluate definite integrals using antiderivatives.
    \item State and apply both parts of the Fundamental Theorem of Calculus.
    \item Translate a rate function into a total accumulated change.
    \item Determine area between a curve and an axis or between two curves when appropriate.
    \item Interpret integrals with correct units.
\end{enumerate}

\subsubsection{Applied Contexts}

\begin{itemize}
    \item Total population change from a population-growth rate.
    \item Total drug exposure from a concentration-rate model.
    \item Accumulated resource consumption.
    \item Net change in economic or demographic variables.
    \item Area-based measures in biological or physical models.
\end{itemize}

\subsubsection{Difficult Topics}

\begin{itemize}
    \item Understanding the difference between signed net accumulation and geometric area.
    \item Moving between a rate function and an accumulated quantity.
    \item Setting up an appropriate definite integral from a verbal context.
    \item Applying the Fundamental Theorem correctly when variable limits are present.
    \item Interpreting units of integrals, especially when integrating a rate.
    \item Recognizing when a graph lies below the horizontal axis.
\end{itemize}

\subsection{Module 6: Techniques and Applications of Integration}

\textbf{Estimated duration:} 1+ week

\subsubsection{Syllabus Chapters and Sections}

\begin{itemize}
    \item Section 7.1: The Substitution Rule
    \item Section 7.2: Integration by Parts and Practicing Integration
    \item Section 7.3: Rational Functions and Partial Fractions
\end{itemize}

\subsubsection{Major Concepts}

\begin{itemize}
    \item Reverse chain rule and substitution.
    \item Change of variables in definite integrals.
    \item Integration by parts:
    \[
        \int u\,dv = uv-\int v\,du.
    \]
    \item Algebraic preparation of integrals.
    \item Rational functions.
    \item Polynomial long division when necessary.
    \item Partial-fraction decomposition.
    \item Selection of an efficient integration technique.
\end{itemize}

\subsubsection{Learning Objectives}

Students will be able to:

\begin{enumerate}
    \item Identify integrals suitable for $u$-substitution.
    \item Select a substitution and transform the integral correctly.
    \item Evaluate definite integrals using substitution and appropriately change the bounds or back-substitute.
    \item Apply integration by parts to products of functions.
    \item Choose useful assignments for $u$ and $dv$.
    \item Simplify integrands algebraically before integrating when appropriate.
    \item Divide polynomials when the degree of the numerator is greater than or equal to the degree of the denominator.
    \item Decompose rational functions into partial fractions.
    \item Integrate rational functions after partial-fraction decomposition.
    \item Verify antiderivatives through differentiation.
\end{enumerate}

\subsubsection{Difficult Topics}

\begin{itemize}
    \item Recognizing a hidden composite-function structure for substitution.
    \item Managing substitutions in definite integrals.
    \item Choosing $u$ and $dv$ strategically in integration by parts.
    \item Repeating integration by parts when necessary.
    \item Factoring denominators correctly for partial fractions.
    \item Determining the correct partial-fraction form for repeated or irreducible factors.
    \item Avoiding algebra mistakes in decomposition coefficients.
    \item Choosing among substitution, integration by parts, partial fractions, and direct integration.
\end{itemize}

\section{Chapter and Topic Map}

\begin{longtable}{p{0.17\textwidth} p{0.15\textwidth} p{0.30\textwidth} p{0.30\textwidth}}
\toprule
\textbf{Chapter/Section} & \textbf{Module} & \textbf{Topic} & \textbf{Primary Skill Outcome} \\
\midrule
\endhead

Chapter 1 & 1 & Basic properties of functions & Analyze functions, domains, ranges, graphs, and transformations. \\

2.1 & 1 & Exponential growth and decay & Build and interpret exponential models. \\

2.2 & 1 & Sequences & Represent and analyze discrete numerical patterns. \\

3.1 & 2 & Limits & Evaluate and interpret limiting behavior. \\

3.2 & 2 & Continuity & Determine continuity and classify discontinuities. \\

3.3 & 2 & Limits at infinity & Describe end behavior and asymptotes. \\

3.4 & 2 & Trigonometric limits and Sandwich Theorem & Establish key trigonometric limits rigorously. \\

3.5 & 2 & Properties of continuous functions & Apply continuity theorems to establish existence. \\

3.6 & 2 & Formal limit definitions & Express limiting behavior using formal mathematical language. \\

4.1 & 3 & Formal derivative definition & Compute and interpret derivatives from first principles. \\

4.2 & 3 & Basic derivative rules & Differentiate sums, constants, and simple functions. \\

4.3 & 3 & Power rule and rules of differentiation & Differentiate polynomial and related functions efficiently. \\

4.4 & 3 & Product and quotient rules & Differentiate products and ratios of functions. \\

4.5 & 3 & Chain rule & Differentiate composite functions. \\

4.6 & 3 & Implicit differentiation & Differentiate relationships not explicitly solved for $y$. \\

4.7 & 3 & Higher derivatives & Analyze acceleration, curvature, and changing rates. \\

4.8 & 3 & Trigonometric derivatives & Differentiate trigonometric models. \\

4.9 & 3 & Exponential derivatives & Analyze continuously growing or decaying quantities. \\

4.10 & 3 & Inverse-function derivatives & Differentiate logarithmic and inverse functions. \\

4.11 & 3 & Linear approximation & Approximate nonlinear functions locally. \\

5.1 & 4 & Extreme and Mean Value Theorems & Justify and locate extrema; analyze average versus instantaneous change. \\

5.2 & 4 & Monotonicity and concavity & Determine increasing/decreasing behavior and graph shape. \\

5.3 & 4 & Extrema and inflection points & Classify local and global behavior. \\

5.4 & 4 & Optimization & Solve constrained maximum/minimum problems. \\

5.5 & 4 & L'Hopital's Rule & Evaluate selected indeterminate limits. \\

5.10 & 5 & Antiderivatives & Reverse differentiation and solve initial-value problems. \\

6.1 & 5 & Definite integral & Quantify net accumulation and signed area. \\

6.2 & 5 & Fundamental Theorem of Calculus & Connect derivatives and definite integrals. \\

6.3 & 5 & Applications of integration & Model accumulated quantities in applied settings. \\

7.1 & 6 & Substitution rule & Evaluate composite-function integrals. \\

7.2 & 6 & Integration by parts & Integrate products and strategically chosen forms. \\

7.3 & 6 & Rational functions and partial fractions & Integrate rational expressions systematically. \\
\bottomrule
\end{longtable}

\section{Most Challenging Topics and Support Plan}

\subsection{High-Difficulty Topics}

\begin{longtable}{p{0.28\textwidth} p{0.35\textwidth} p{0.28\textwidth}}
\toprule
\textbf{Topic} & \textbf{Why It Is Challenging} & \textbf{Recommended Preparation} \\
\midrule
\endhead

Formal limits and $\varepsilon$--$\delta$ definitions &
Requires abstract reasoning about arbitrary closeness and logical quantifiers. &
Master graphical and numerical limits before formal proofs. \\

Derivative from first principles &
Combines limit reasoning, algebraic expansion, simplification, and conceptual interpretation. &
Review factoring, rationalization, and the difference quotient. \\

Chain rule &
Students must recognize nested functional structure rather than merely apply memorized formulas. &
Practice identifying ``outer'' and ``inner'' functions. \\

Implicit differentiation &
Requires treating $y$ as a function of $x$ even when it is not explicitly isolated. &
Review chain rule and derivative notation such as $\frac{dy}{dx}$. \\

Optimization &
Requires translating prose into equations, selecting a valid domain, differentiating, and interpreting a result. &
Practice function modeling and unit analysis. \\

Fundamental Theorem of Calculus &
Connects two major ideas---rates of change and accumulation---that can initially seem unrelated. &
Develop strong understanding of antiderivatives and definite integrals. \\

Integration by parts &
Requires strategic selection of $u$ and $dv$ rather than a single mechanical rule. &
Review product rule and common derivative--antiderivative pairs. \\

Partial fractions &
Involves factoring, algebraic decomposition, systems of equations, and integration. &
Review polynomial factoring and rational expressions. \\
\bottomrule
\end{longtable}

\subsection{Suggested Study Sequence}

\begin{enumerate}
    \item Build algebra and function fluency before attempting calculus procedures.
    \item Treat limits as the conceptual foundation for derivatives.
    \item Learn derivative rules only after understanding the derivative as a limit and a rate of change.
    \item Practice derivative applications after computational differentiation becomes reliable.
    \item Learn integration as accumulation and reversal of differentiation before focusing on advanced integration techniques.
    \item Use a technique-selection checklist for integration:
    \[
    \text{Direct rule} \rightarrow
    \text{Substitution} \rightarrow
    \text{Integration by parts} \rightarrow
    \text{Partial fractions}.
    \]
\end{enumerate}

\section{Conceptual Dependency Map}

\[
\boxed{\text{Algebra, Functions, and Graphs}}
\]

\[
\downarrow
\]

\[
\boxed{\text{Exponential Models and Sequences}}
\]

\[
\downarrow
\]

\[
\boxed{\text{Limits and Continuity}}
\]

\[
\downarrow
\]

\[
\boxed{\text{Derivative Definition}}
\]

\[
\downarrow
\]

\[
\boxed{\text{Derivative Rules}}
\]

\[
\downarrow
\]

\[
\boxed{\text{Applications of Derivatives}}
\]

\[
\downarrow
\]

\[
\boxed{\text{Antiderivatives and Definite Integrals}}
\]

\[
\downarrow
\]

\[
\boxed{\text{Fundamental Theorem of Calculus}}
\]

\[
\downarrow
\]

\[
\boxed{\text{Applications and Techniques of Integration}}
\]

\section{Assessment-Aligned Competencies}

A student who has successfully completed this course should be prepared to:

\begin{itemize}
    \item Solve routine and non-routine calculus problems symbolically.
    \item Interpret graphical and numerical information in calculus contexts.
    \item Model biological and social-science phenomena using functions, rates, and accumulated quantities.
    \item Explain the practical meaning of a derivative or integral, including its units.
    \item Select appropriate calculus tools instead of applying procedures mechanically.
    \item Check results for reasonableness using graphs, estimates, units, endpoint behavior, or differentiation.
    \item Transition to subsequent coursework involving multivariable calculus, differential equations, statistics, mathematical biology, economics, or quantitative social-science modeling.
\end{itemize}

\end{document}